% =========================================================
% Subsets of R --- Notes
% Chapter: functions
% Section: subsets-real-line
% Volume III $\cdot$ Analysis
% =========================================================

\subsection*{\texorpdfstring{Subsets of $\mathbb{R}$}{Subsets of R}}
\label{sec:subsets-of-R}
\begin{remark*}[Interpretation]
Cluster point = limit point = accumulation point. Bartle and Lebl use
\emph{cluster point}; Rudin uses \emph{limit point}; Tao uses both.
Adherent point (Tao) is the weaker notion that includes isolated points.
\end{remark*}

% ---------------------------------------------------------
% Definition --- Centered Open Interval
% ---------------------------------------------------------

\begin{definitionbox}{Definition (Centered Open Interval)}
\begin{definition}[Centered Open Interval]
\label{def:epsilon-neighbourhood}
Let $x\in\mathbb{R}$ and let $\varepsilon>0$. The
\textbf{centered open interval of radius $\varepsilon$ about $x$} is
\[
  I_\varepsilon(x):=(x-\varepsilon,x+\varepsilon)
  =\{y\in\mathbb{R}: |y-x|<\varepsilon\}.
\]
Older sources often call this the $\varepsilon$-neighbourhood of $x$ and write
$V_\varepsilon(x)$.
\end{definition}
\end{definitionbox}

\begin{remark*}[Standard quantified statement]
\[
  y\in I_\varepsilon(x)\;\Longleftrightarrow\; |y-x|<\varepsilon.
\]
\end{remark*}

\begin{remark*}[Interpretation]
The centered open interval is the concrete real-line local region used before
the later topological language of neighbourhoods.
\end{remark*}

\begin{dependencies}
\begin{itemize}
  \item \hyperref[def:open-ball]{Open Ball on the Real Line}
\end{itemize}
\end{dependencies}

\begin{definitionbox}{Definition (Punctured Centered Open Interval)}
\begin{definition}[Punctured Centered Open Interval]
\label{def:deleted-epsilon-neighbourhood}
Let $x\in\mathbb{R}$ and let $\varepsilon>0$. The
\textbf{punctured centered open interval of radius $\varepsilon$ about $x$} is
\[
  I_\varepsilon^\times(x):=I_\varepsilon(x)\setminus\{x\}
  =\{y\in\mathbb{R}:0<|y-x|<\varepsilon\}.
\]
Older sources often call this the deleted $\varepsilon$-neighbourhood of $x$
and write $V_\varepsilon^*(x)$.
\end{definition}
\end{definitionbox}

\begin{remark*}[Standard quantified statement]
\[
  y\in I_\varepsilon^\times(x)\;\Longleftrightarrow\;0<|y-x|<\varepsilon.
\]
\end{remark*}

\begin{remark*}[Interpretation]
The punctured centered open interval removes the base point. It is the concrete
input region used in limit definitions.
\end{remark*}

\begin{dependencies}
\begin{itemize}
  \item \hyperref[def:epsilon-neighbourhood]{Centered Open Interval}
\end{itemize}
\end{dependencies}

% ---------------------------------------------------------
% Definition --- Cluster Point
% ---------------------------------------------------------

\begin{definitionbox}{Definition (Cluster Point)}
\begin{definition}[Cluster Point]
\label{def:cluster-point-r}
Let $X \subseteq \mathbb{R}$ and $x \in \mathbb{R}$. The point $x$
is a \textbf{cluster point} of $X$ if
\[
  \forall\varepsilon > 0,\;
  \exists\, y \in X \setminus \{x\} :\;
  |y-x|<\varepsilon.
\]
\end{definition}
\end{definitionbox}

\begin{remark*}[Standard quantified statement]
\[
  x\text{ is a cluster point of }X
  \;\iff\;
  \forall\varepsilon > 0,\;
  \exists\, y \in X \setminus \{x\} :\;
  |y - x| < \varepsilon.
\]
\end{remark*}

\begin{remark*}[Predicate reading]
\[
  x\text{ is a cluster point of }X
  \;\colonequiv\;
  \forall\varepsilon > 0,\;
  I_\varepsilon^\times(x) \cap X \neq \emptyset.
\]
\end{remark*}

\begin{remark*}[Negated quantified statement]
\[
  \neg\,x\text{ is a cluster point of }X
  \;\iff\;
  \exists\varepsilon > 0,\;
  \forall\, y \in X \setminus \{x\} :\;
  \neg|y-x|<\varepsilon.
\]
\end{remark*}

\begin{remark*}[Negation predicate reading]
\[
  \neg\,x\text{ is a cluster point of }X
  \;\colonequiv\;
  \exists\varepsilon > 0 :\;
  I_\varepsilon^\times(x) \cap X = \emptyset.
\]
\end{remark*}

\begin{remark*}[Interpretation]
A cluster point is approached by $X$ via distinct witnesses
arbitrarily close. A cluster point need not belong to $X$: $0$ is
a cluster point of $(0,1)$ but $0 \notin (0,1)$. The condition
$\lim_{x\to c}f(x)=L$ requires
$\operatorname{IsCluster}(c, \operatorname{dom}(f))$; without
it the punctured ball condition is vacuous.
\end{remark*}

\begin{dependencies}
\begin{itemize}
  \item \hyperref[def:deleted-epsilon-neighbourhood]{Punctured Centered Open Interval}
\end{itemize}
\end{dependencies}

% ---------------------------------------------------------
% Theorem --- Sequential Characterization of Cluster Points
% ---------------------------------------------------------

\begin{theorembox}{Theorem}
\begin{theorem}[Sequential Characterization of Cluster Points]
\label{thm:cluster-point-sequential}
$c\text{ is a cluster point of }A$ if and only if there exists
$(a_n) \subseteq A \setminus \{c\}$ with $a_n \to c$.

\smallskip\noindent
\hyperref[prf:cluster-point-sequential]{\textit{Go to proof.}}
\end{theorem}
\end{theorembox}

\begin{remark*}[Standard quantified statement]
\[
  c\text{ is a cluster point of }A
  \;\iff\;
  \exists (a_n) \subseteq A \setminus \{c\} :\;
  a_n\to c.
\]
\end{remark*}

\begin{remark*}[Interpretation]
The forward direction: use $\varepsilon = 1/n$ to build the witness
sequence. The reverse direction is immediate since the sequence
witnesses the ball condition for every $\varepsilon$. This bridges
the $\varepsilon$-$\delta$ definition with the sequential language.
\end{remark*}

\begin{dependencies}
\begin{itemize}
  \item \hyperref[def:cluster-point-r]{Cluster Point}
\end{itemize}
\end{dependencies}

% ---------------------------------------------------------
% Definition --- Adherent Point and Isolated Point
% ---------------------------------------------------------

\begin{definitionbox}{Definition (Adherent Point)}
\begin{definition}[Adherent Point]
\label{def:adherent-point-r}
Let $X\subseteq\mathbb{R}$ and let $x\in\mathbb{R}$. The point $x$ is
an \textbf{adherent point} of $X$ if
\[
  \forall\varepsilon>0\;\exists y\in X\;(|y-x|<\varepsilon).
\]
\end{definition}
\end{definitionbox}

\begin{remark*}[Standard quantified statement]
\[
  \forall\varepsilon>0\;\bigl(I_\varepsilon(x)\cap X\neq\emptyset\bigr).
\]
\end{remark*}

\begin{remark*}[Predicate reading]
\[
  \operatorname{IsAdherent}(x,X)
  \;\coloneqq\;
  \forall\varepsilon>0\;\exists y\in X\;(|y-x|<\varepsilon).
\]
\end{remark*}

\begin{remark*}[Negated quantified statement]
\[
  \exists\varepsilon>0\;\forall y\in X\;(|y-x|\geq\varepsilon).
\]
\end{remark*}

\begin{remark*}[Interpretation]
An adherent point can be witnessed by points of $X$ arbitrarily close to
$x$, including possibly $x$ itself.
\end{remark*}

\begin{dependencies}
\begin{itemize}
  \item \hyperref[def:epsilon-neighbourhood]{Centered Open Interval}
\end{itemize}
\end{dependencies}

\begin{definitionbox}{Definition (Isolated Point)}
\begin{definition}[Isolated Point]
\label{def:isolated-point-r}
Let $X\subseteq\mathbb{R}$ and let $x\in X$. The point $x$ is an
\textbf{isolated point} of $X$ if
\[
  \exists\varepsilon>0\;\bigl(I_\varepsilon(x)\cap X=\{x\}\bigr).
\]
\end{definition}
\end{definitionbox}

\begin{remark*}[Standard quantified statement]
\[
  \exists\varepsilon>0\;\forall y\in X\;\bigl(|y-x|<\varepsilon\Rightarrow y=x\bigr).
\]
\end{remark*}

\begin{remark*}[Predicate reading]
\[
  \operatorname{IsIsolated}(x,X)
  \;\coloneqq\;
  x\in X\wedge\exists\varepsilon>0\;\bigl(I_\varepsilon(x)\cap X=\{x\}\bigr).
\]
\end{remark*}

\begin{remark*}[Negated quantified statement]
\[
  \forall\varepsilon>0\;\exists y\in X\setminus\{x\}\;(|y-x|<\varepsilon).
\]
\end{remark*}

\begin{remark*}[Interpretation]
An isolated point belongs to the set but has a centered open interval containing
no other points of the set.
\end{remark*}

\begin{dependencies}
\begin{itemize}
  \item \hyperref[def:adherent-point-r]{Adherent Point}
  \item \hyperref[def:cluster-point-r]{Cluster Point}
\end{itemize}
\end{dependencies}

% ---------------------------------------------------------
% Definition --- Interior, Exterior, Boundary
% ---------------------------------------------------------

\begin{definitionbox}{Definition (Interior Point)}
\begin{definition}[Interior Point]
\label{def:interior-point-r}
Let $X\subseteq\mathbb{R}$ and let $x\in\mathbb{R}$. The point $x$ is
an \textbf{interior point} of $X$ if
\[
  \exists\varepsilon>0\;\bigl(I_\varepsilon(x)\subseteq X\bigr).
\]
\end{definition}
\end{definitionbox}

\begin{remark*}[Standard quantified statement]
\[
  \exists\varepsilon>0\;\forall y\in\mathbb{R}\;\bigl(|y-x|<\varepsilon\Rightarrow y\in X\bigr).
\]
\end{remark*}

\begin{remark*}[Predicate reading]
\[
  \operatorname{IsInterior}(x,X)\;\coloneqq\;\exists\varepsilon>0\;\bigl(I_\varepsilon(x)\subseteq X\bigr).
\]
\end{remark*}

\begin{remark*}[Interpretation]
An interior point has a full centered open interval contained in the set.
\end{remark*}

\begin{dependencies}
\begin{itemize}
  \item \hyperref[def:epsilon-neighbourhood]{Centered Open Interval}
\end{itemize}
\end{dependencies}

\begin{definitionbox}{Definition (Boundary Point)}
\begin{definition}[Boundary Point]
\label{def:boundary-point-r}
Let $X\subseteq\mathbb{R}$ and let $x\in\mathbb{R}$. The point $x$ is
a \textbf{boundary point} of $X$ if
\[
  \forall\varepsilon>0\;\bigl(I_\varepsilon(x)\cap X\neq\emptyset\bigr)
  \wedge\bigl(I_\varepsilon(x)\cap X^c\neq\emptyset\bigr).
\]
\end{definition}
\end{definitionbox}

\begin{remark*}[Standard quantified statement]
\[
  \forall\varepsilon>0\;\bigl(I_\varepsilon(x)\cap X\neq\emptyset\bigr)
  \wedge\bigl(I_\varepsilon(x)\cap X^c\neq\emptyset\bigr).
\]
\end{remark*}

\begin{remark*}[Predicate reading]
\[
  \operatorname{IsOnBoundary}(x,X)\;\coloneqq\;
  \forall\varepsilon>0\;\bigl(I_\varepsilon(x)\cap X\neq\emptyset\bigr)
  \wedge\bigl(I_\varepsilon(x)\cap X^c\neq\emptyset\bigr).
\]
\end{remark*}

\begin{remark*}[Interpretation]
A boundary point sees both the set and its complement in every
centered open interval.
\end{remark*}

\begin{dependencies}
\begin{itemize}
  \item \hyperref[def:epsilon-neighbourhood]{Centered Open Interval}
\end{itemize}
\end{dependencies}

\begin{definitionbox}{Definition (Interior)}
\begin{definition}[Interior]
\label{def:interior-r}
Let $X\subseteq\mathbb{R}$. The \textbf{interior} of $X$ is
\[
  X^\circ:=\{x\in\mathbb{R}:\exists\varepsilon>0\;I_\varepsilon(x)\subseteq X\}.
\]
\end{definition}
\end{definitionbox}

\begin{remark*}[Standard quantified statement]
\[
  x\in X^\circ\;\Longleftrightarrow\;\exists\varepsilon>0\;I_\varepsilon(x)\subseteq X.
\]
\end{remark*}

\begin{remark*}[Interpretation]
The interior collects exactly the points around which the set contains a
whole centered open interval.
\end{remark*}

\begin{dependencies}
\begin{itemize}
  \item \hyperref[def:interior-point-r]{Interior Point}
\end{itemize}
\end{dependencies}

\begin{definitionbox}{Definition (Boundary)}
\begin{definition}[Boundary]
\label{def:boundary-r}
Let $X\subseteq\mathbb{R}$. The \textbf{boundary} of $X$ is
\[
  \partial X:=\{x\in\mathbb{R}:\forall\varepsilon>0\;(I_\varepsilon(x)\cap X\neq\emptyset\wedge I_\varepsilon(x)\cap X^c\neq\emptyset)\}.
\]
\end{definition}
\end{definitionbox}

\begin{remark*}[Standard quantified statement]
\[
  x\in\partial X\;\Longleftrightarrow\;
  \forall\varepsilon>0\;(I_\varepsilon(x)\cap X\neq\emptyset\wedge I_\varepsilon(x)\cap X^c\neq\emptyset).
\]
\end{remark*}

\begin{remark*}[Interpretation]
The boundary records the points where every centered open interval touches both
inside and outside of the set.
\end{remark*}

\begin{dependencies}
\begin{itemize}
  \item \hyperref[def:boundary-point-r]{Boundary Point}
\end{itemize}
\end{dependencies}

% ---------------------------------------------------------
% Definition --- Closure
% ---------------------------------------------------------
\begin{definitionbox}{Definition (Closure)}
\begin{definition}[Closure]
\label{def:closure-r}
\[
\overline{X}
:=
\{x\in\mathbb{R}:\forall\varepsilon>0\;I_\varepsilon(x)\cap X\neq\emptyset\}.
\]
\end{definition}
\end{definitionbox}

\begin{remark*}[Standard quantified statement]
\[
  x \in \overline{X}
  \;\iff\;
  x\text{ is an adherent point of }X
  \;\iff\;
  \forall\varepsilon>0 : I_\varepsilon(x)\cap X\neq\emptyset.
\]
\end{remark*}

\begin{remark*}[Interpretation]
The closure adds all adherent points of $X$, so it fills in the limiting
points that can be reached by points of $X$.


$\overline{(a,b)} = [a,b]$;\quad $\overline{\mathbb{Q}} = \mathbb{R}$;\quad
$\overline{\mathbb{Z}} = \mathbb{Z}$;\quad
$\overline{\{1/n\}} = \{1/n\}\cup\{0\}$.
\end{remark*}

\begin{dependencies}
\begin{itemize}
  \item \hyperref[def:adherent-point-r]{Adherent Point}
\end{itemize}
\end{dependencies}

\begin{propositionbox}{Proposition}
\begin{proposition}[Adherent Points Are Exactly Closure Points]
\label{prop:adherent-points-are-closure-points}
Let \(X\subseteq\mathbb{R}\) and \(x\in\mathbb{R}\).  Then
\[
  x\in\overline{X}
  \quad\Longleftrightarrow\quad
  x\text{ is an adherent point of }X.
\]

\smallskip\noindent
\hyperref[prf:adherent-points-are-closure-points]{\textit{Go to proof.}}
\end{proposition}
\end{propositionbox}

\begin{remark*}[Standard quantified statement]
\[
  x\in\overline{X}
  \;\Longleftrightarrow\;
  \operatorname{IsAdherent}(x,X).
\]
\end{remark*}

\begin{remark*}[Interpretation]
Closure is the set-level packaging of the pointwise adherent-point predicate.
\end{remark*}

\begin{dependencies}
\begin{itemize}
  \item \hyperref[def:adherent-point-r]{Adherent Point}
  \item \hyperref[def:closure-r]{Closure}
\end{itemize}
\end{dependencies}

\begin{propositionbox}{Proposition}
\begin{proposition}[Isolated Points Are Noncluster Adherent Points]
\label{prop:isolated-points-are-noncluster-adherent-points}
Let \(X\subseteq\mathbb{R}\) and \(x\in X\).  Then \(x\) is isolated in
\(X\) if and only if \(x\) is an adherent point of \(X\) and \(x\) is not
a cluster point of \(X\).

\smallskip\noindent
\hyperref[prf:isolated-points-are-noncluster-adherent-points]{\textit{Go to proof.}}
\end{proposition}
\end{propositionbox}

\begin{remark*}[Standard quantified statement]
\[
  x\in X
  \;\Longrightarrow\;
  \bigl(\operatorname{IsIsolated}(x,X)
  \Longleftrightarrow
  \operatorname{IsAdherent}(x,X)\wedge\neg\operatorname{IsCluster}(x,X)\bigr).
\]
\end{remark*}

\begin{remark*}[Interpretation]
The point itself witnesses adherence, while the absence of distinct nearby
points prevents cluster behavior.
\end{remark*}

\begin{dependencies}
\begin{itemize}
  \item \hyperref[def:isolated-point-r]{Isolated Point}
  \item \hyperref[def:adherent-point-r]{Adherent Point}
  \item \hyperref[def:cluster-point-r]{Cluster Point}
\end{itemize}
\end{dependencies}

\begin{propositionbox}{Proposition}
\begin{proposition}[Interior Membership Characterization]
\label{prop:interior-membership-characterization}
Let \(X\subseteq\mathbb{R}\) and \(x\in\mathbb{R}\).  Then
\[
  x\in X^\circ
  \quad\Longleftrightarrow\quad
  x\text{ is an interior point of }X.
\]

\smallskip\noindent
\hyperref[prf:interior-membership-characterization]{\textit{Go to proof.}}
\end{proposition}
\end{propositionbox}

\begin{remark*}[Standard quantified statement]
\[
  x\in X^\circ
  \;\Longleftrightarrow\;
  x\text{ is an interior point of }X.
\]
\end{remark*}

\begin{remark*}[Interpretation]
The interior set collects exactly the points satisfying the interior-point
predicate.
\end{remark*}

\begin{dependencies}
\begin{itemize}
  \item \hyperref[def:interior-r]{Interior}
  \item \hyperref[def:interior-point-r]{Interior Point}
\end{itemize}
\end{dependencies}

\begin{propositionbox}{Proposition}
\begin{proposition}[Interior Is Contained in the Set]
\label{prop:interior-is-contained-in-set}
For every \(X\subseteq\mathbb{R}\), \(X^\circ\subseteq X\).

\smallskip\noindent
\hyperref[prf:interior-is-contained-in-set]{\textit{Go to proof.}}
\end{proposition}
\end{propositionbox}

\begin{remark*}[Standard quantified statement]
\[
  X^\circ\subseteq X.
\]
\end{remark*}

\begin{remark*}[Interpretation]
Interior points have a whole neighborhood inside \(X\), so in particular the
center point lies in \(X\).
\end{remark*}

\begin{dependencies}
\begin{itemize}
  \item \hyperref[def:interior-r]{Interior}
  \item \hyperref[def:epsilon-neighbourhood]{Centered Open Interval}
\end{itemize}
\end{dependencies}

\begin{propositionbox}{Proposition}
\begin{proposition}[Closure Is the Smallest Closed Superset]
\label{prop:closure-is-smallest-closed-superset}
Let \(X\subseteq\mathbb{R}\).  Then \(X\subseteq\overline{X}\),
\(\overline{X}\) is closed, and if \(C\subseteq\mathbb{R}\) is closed with
\(X\subseteq C\), then \(\overline{X}\subseteq C\).

\smallskip\noindent
\hyperref[prf:closure-is-smallest-closed-superset]{\textit{Go to proof.}}
\end{proposition}
\end{propositionbox}

\begin{remark*}[Standard quantified statement]
\[
  X\subseteq\overline{X}
  \quad\text{and}\quad
  C\text{ closed},\,X\subseteq C\Rightarrow \overline{X}\subseteq C.
\]
\end{remark*}

\begin{remark*}[Interpretation]
The closure is not just a closed superset; it is the minimal closed set that
contains \(X\).
\end{remark*}

\begin{dependencies}
\begin{itemize}
  \item \hyperref[def:closure-r]{Closure}
  \item \hyperref[def:closed-set-r]{Closed Subset of \(\mathbb{R}\)}
  \item \hyperref[lem:closure-elementary]{Elementary Properties of Closures}
\end{itemize}
\end{dependencies}

\begin{propositionbox}{Proposition}
\begin{proposition}[Boundary as Closure Minus Interior]
\label{prop:boundary-as-closure-minus-interior}
For every \(X\subseteq\mathbb{R}\),
\[
  \partial X=\overline{X}\setminus X^\circ.
\]

\smallskip\noindent
\hyperref[prf:boundary-as-closure-minus-interior]{\textit{Go to proof.}}
\end{proposition}
\end{propositionbox}

\begin{remark*}[Standard quantified statement]
\[
  \partial X=\overline{X}\setminus X^\circ.
\]
\end{remark*}

\begin{remark*}[Interpretation]
Boundary points are adherent to \(X\) but do not have a whole neighborhood
lying inside \(X\).
\end{remark*}

\begin{dependencies}
\begin{itemize}
  \item \hyperref[def:boundary-r]{Boundary}
  \item \hyperref[def:closure-r]{Closure}
  \item \hyperref[def:interior-r]{Interior}
\end{itemize}
\end{dependencies}

% ---------------------------------------------------------
% Lemma --- Elementary Properties of Closures
% ---------------------------------------------------------

\begin{lemmabox}{Lemma}
\begin{lemma}[Elementary Properties of Closures]
\label{lem:closure-elementary}
Let $X, Y \subseteq \mathbb{R}$. Then: (i) $X \subseteq \overline{X}$;
(ii) $\overline{X \cup Y} = \overline{X} \cup \overline{Y}$;
(iii) $\overline{X \cap Y} \subseteq \overline{X} \cap \overline{Y}$;
(iv) $X \subseteq Y \Rightarrow \overline{X} \subseteq \overline{Y}$.

\smallskip\noindent
\hyperref[prf:closure-elementary]{\textit{Go to proof.}}
\end{lemma}
\end{lemmabox}

\begin{remark*}[Standard quantified statement]
\[
  \overline{X \cup Y} = \overline{X} \cup \overline{Y};\quad
  \overline{X \cap Y} \subseteq \overline{X} \cap \overline{Y}.
\]
\end{remark*}

\begin{remark*}[Interpretation]
Take $X = (0,1)$, $Y = (1,2)$: $\overline{X\cap Y} = \emptyset$ but
$\overline{X}\cap\overline{Y} = \{1\}$.


Closure preserves inclusions and finite unions, but intersection can lose
limit points that are approached separately by the two sets.
\end{remark*}

\begin{dependencies}
\begin{itemize}
  \item \hyperref[def:closure-r]{Closure}
\end{itemize}
\end{dependencies}

% ---------------------------------------------------------
% Definition --- Closed Set
% ---------------------------------------------------------

\begin{definitionbox}{Definition (Closed Subset of $\mathbb{R}$)}
\begin{definition}[Closed Subset of $\mathbb{R}$]
\label{def:closed-set-r}
$E\text{ is closed} \iff \overline{E} = E$.
\end{definition}
\end{definitionbox}

\begin{remark*}[Standard quantified statement]
\[
  E\text{ is closed}
  \;\iff\;
  \forall x \in \mathbb{R} :\;
  x\text{ is an adherent point of }E \Rightarrow x \in E.
\]
\end{remark*}

\begin{remark*}[Predicate reading]
\[
  E\text{ is closed}
  \;\colonequiv\;
  \forall x :\;
  x\text{ is an adherent point of }E \Rightarrow x \in E.
\]
\end{remark*}

\begin{remark*}[Negated quantified statement]
\[
  \neg\,E\text{ is closed}
  \;\iff\;
  \exists x \notin E :\; x\text{ is an adherent point of }E.
\]
\end{remark*}

\begin{remark*}[Negation predicate reading]
\[
  \neg\,E\text{ is closed}
  \;\colonequiv\;
  \exists x \notin E :\;
  \forall\varepsilon>0 : I_\varepsilon(x)\cap E\neq\emptyset.
\]
\end{remark*}

\begin{remark*}[Interpretation]
$E$ is closed iff it contains its own boundary. Examples:
$[a,b]$ closed; $(a,b)$ not closed; $\mathbb{R}$ and $\emptyset$
both closed; $\mathbb{Q}$ not closed.
\end{remark*}

\begin{dependencies}
\begin{itemize}
  \item \hyperref[def:closure-r]{Closure}
\end{itemize}
\end{dependencies}

% ---------------------------------------------------------
% Corollary --- Sequentially Closed
% ---------------------------------------------------------

\begin{corollarybox}{Corollary}
\begin{corollary}[Closed iff Sequentially Closed]
\label{cor:closed-iff-seq-limits}
$X\text{ is closed} \iff \forall (a_n)\subseteq X,\;
a_n \to x \Rightarrow x \in X$.

\smallskip\noindent
\hyperref[prf:closed-iff-seq-limits]{\textit{Go to proof.}}
\end{corollary}
\end{corollarybox}

\begin{remark*}[Standard quantified statement]
\[
  X\text{ is closed}
  \;\iff\;
  \forall (a_n)\subseteq X,\;
  a_n\to x
  \Rightarrow x \in X.
\]
\end{remark*}

\begin{remark*}[Interpretation]
A closed set cannot be escaped by convergent sequences. Every adherent
point has a witness sequence; closedness forces the limit into the set.
\end{remark*}

\begin{dependencies}
\begin{itemize}
  \item \hyperref[def:closed-set-r]{Closed Subset of $\mathbb{R}$}
  \item \hyperref[thm:cluster-point-sequential]{Sequential Characterization of Cluster Points}
\end{itemize}
\end{dependencies}

% ---------------------------------------------------------
% Corollary --- Every Point of an Interval Is a Cluster Point
% ---------------------------------------------------------

\begin{corollarybox}{Corollary}
\begin{corollary}[Every Point of an Interval Is a Cluster Point]
\label{cor:interval-all-limit-points}
If $I$ is an interval, then every $x\in I$ is a cluster point of $I$.

\smallskip\noindent
\hyperref[prf:interval-all-limit-points]{\textit{Go to proof.}}
\end{corollary}
\end{corollarybox}

\begin{remark*}[Standard quantified statement]
\[
  \forall x\in I\;\forall\varepsilon>0\;\exists y\in I\setminus\{x\}\;(|y-x|<\varepsilon).
\]
\end{remark*}

\begin{remark*}[Predicate reading]
The displayed quantified statement is the predicate-level form of $label: its hypotheses name the input conditions, and its conclusion names the asserted structure or relation.
\end{remark*}

\begin{remark*}[Interpretation]
This justifies intervals as natural domains for function limits.
$\lim_{x\to c}f(x)=L$ requires
$\operatorname{IsCluster}(c, \operatorname{dom}(f))$; on any
interval this is free.
\end{remark*}

\begin{dependencies}
\begin{itemize}
  \item \hyperref[def:cluster-point-r]{Cluster Point}
\end{itemize}
\end{dependencies}

% ---------------------------------------------------------
% Definition --- Bounded Subset
% ---------------------------------------------------------

\begin{definitionbox}{Definition (Bounded Subset of $\mathbb{R}$)}
\begin{definition}[Bounded Subset of $\mathbb{R}$]
\label{def:bounded-set-r}
$X\text{ is bounded} \iff \exists M>0,\;\forall x\in X:
|x|\leq M$.
\end{definition}
\end{definitionbox}

\begin{remark*}[Standard quantified statement]
\[
  \exists M>0\;\forall x\in X\;(|x|\leq M).
\]
\end{remark*}

\begin{remark*}[Predicate reading]
The displayed quantified statement is the predicate-level form of $label: its hypotheses name the input conditions, and its conclusion names the asserted structure or relation.
\end{remark*}

\begin{remark*}[Negated quantified statement]
\[
  \forall M>0\;\exists x\in X\;(|x|>M).
\]
\end{remark*}

\begin{remark*}[Interpretation]
A bounded subset of $\mathbb{R}$ fits inside some symmetric interval
$[-M,M]$.
\end{remark*}

\begin{dependencies}
\begin{itemize}
  \item \hyperref[def:subset]{Subset}
  \item \hyperref[def:real-chapter-absolute-value-on-r]{Absolute Value on $\mathbb{R}$}
  \item \hyperref[def:order-on-r]{Order on $\mathbb{R}$}
\end{itemize}
\end{dependencies}

% ---------------------------------------------------------
% Theorem --- Heine-Borel
% ---------------------------------------------------------

\begin{theorembox}{Theorem}
\begin{theorem}[Heine--Borel Theorem for $\mathbb{R}$]
\label{thm:heine-borel-subsets-real-line}
$X\text{ is closed} \wedge X\text{ is bounded}$
$\iff$ every $(a_n)\subseteq X$ has a subsequence converging to
a limit in $X$.

\smallskip\noindent
\hyperref[prf:heine-borel-subsets-real-line]{\textit{Go to proof.}}
\end{theorem}
\end{theorembox}

\begin{remark*}[Standard quantified statement]
\[
  X\text{ is closed} \wedge X\text{ is bounded}
  \;\iff\;
  \forall (a_n)\subseteq X,\;
  \exists (a_{n_k}),\;\exists L\in X : a_{n_k}\to L.
\]
\end{remark*}

\begin{remark*}[Predicate reading]
\[
  \operatorname{IsSequentiallyCompact}(X)
  \;\colonequiv\;
  X\text{ is closed} \wedge X\text{ is bounded}.
\]
\end{remark*}

\begin{remark*}[Negated quantified statement]
\[
  \neg\bigl(X\text{ is closed}\wedge X\text{ is bounded}\bigr)
  \;\iff\;
  \exists (a_n)\subseteq X:\;
  \text{no subsequence converges to a limit in } X.
\]
\end{remark*}

\begin{remark*}[Failure modes]
\begin{description}
  \item[Exposition.]
The compactness conclusion fails when either closedness or boundedness fails.
\end{description}
\end{remark*}

\begin{remark*}[Failure modes]
\begin{description}
  \item[Exposition.]
Either $X\text{ is bounded}$ fails --- a sequence escapes to
$\pm\infty$ --- or $X\text{ is closed}$ fails --- a convergent
subsequence exists but its limit lies outside $X$.
\end{description}
\end{remark*}

\begin{remark*}[Interpretation]
Closed prevents limits from escaping through boundary gaps; bounded
prevents escape to infinity. Both are required. Forward: bounded
gives Bolzano--Weierstrass; closed keeps the limit in $X$. Reverse:
negate either hypothesis to construct a sequence with no convergent
subsequence in $X$.
\end{remark*}

\begin{dependencies}
\begin{itemize}
  \item \hyperref[def:closed-set-r]{Closed Subset of $\mathbb{R}$}
  \item \hyperref[def:bounded-set-r]{Bounded Subset of $\mathbb{R}$}
\end{itemize}
\end{dependencies}

% ---------------------------------------------------------
% Definition --- True Near a Point
% ---------------------------------------------------------

\begin{definitionbox}{Definition (True Near a Point)}
\begin{definition}[True Near a Point]
\label{def:true-near}
Property $P(f(x))$ is \textbf{true near $x_0$} if
$\exists\delta>0,\;\forall x : 0<|x-x_0|<
\delta) \Rightarrow P(f(x))$.
\end{definition}
\end{definitionbox}

\begin{remark*}[Standard quantified statement]
\[
  P(f(x))\text{ true near }x_0
  \;\iff\;
  \exists\delta>0,\;\forall x :
  0 < |x-x_0| < \delta \Rightarrow P(f(x)).
\]
\end{remark*}

\begin{remark*}[Interpretation]
$\lim_{x\to x_0}f(x)=L$ states: for every
$\varepsilon>0$, the property $|f(x)-L|<\varepsilon$ is true near $x_0$.
\end{remark*}

\begin{dependencies}
\begin{itemize}
  \item \hyperref[def:epsilon-neighbourhood]{Centered Open Interval}
\end{itemize}
\end{dependencies}

\begin{propositionbox}{Proposition}
\begin{proposition}[True Near Is Stable Under Shrinking]
\label{prop:true-near-stable-under-shrinking}
If a property \(P(f(x))\) is true near \(x_0\), then it remains true on
every sufficiently smaller deleted neighborhood of \(x_0\).

\smallskip\noindent
\hyperref[prf:true-near-stable-under-shrinking]{\textit{Go to proof.}}
\end{proposition}
\end{propositionbox}

\begin{remark*}[Standard quantified statement]
\[
  P(f(x))\text{ true near }x_0
  \;\Longrightarrow\;
  \exists\eta>0\;\forall x\;(0<|x-x_0|<\eta\Rightarrow P(f(x))).
\]
\end{remark*}

\begin{remark*}[Interpretation]
Once a local property holds on one punctured neighborhood, it also holds after
the radius is decreased.
\end{remark*}

\begin{dependencies}
\begin{itemize}
  \item \hyperref[def:true-near]{True Near a Point}
  \item \hyperref[def:deleted-epsilon-neighbourhood]{Punctured Centered Open Interval}
\end{itemize}
\end{dependencies}

\begin{propositionbox}{Proposition}
\begin{proposition}[True Near Is Stable Under Conjunction]
\label{prop:true-near-stable-under-conjunction}
If \(P(f(x))\) is true near \(x_0\) and \(Q(f(x))\) is true near \(x_0\),
then \(P(f(x))\wedge Q(f(x))\) is true near \(x_0\).

\smallskip\noindent
\hyperref[prf:true-near-stable-under-conjunction]{\textit{Go to proof.}}
\end{proposition}
\end{propositionbox}

\begin{remark*}[Standard quantified statement]
\[
  P(f(x))\text{ true near }x_0
  \;\wedge\;
  Q(f(x))\text{ true near }x_0
  \;\Longrightarrow\;
  P(f(x))\wedge Q(f(x))\text{ true near }x_0.
\]
\end{remark*}

\begin{remark*}[Interpretation]
Two local facts can be made simultaneous by taking the smaller of their two
radii.
\end{remark*}

\begin{dependencies}
\begin{itemize}
  \item \hyperref[def:true-near]{True Near a Point}
  \item \hyperref[lem:minimum-of-positive-numbers-is-positive]{Minimum of Positive Numbers Is Positive}
\end{itemize}
\end{dependencies}
